\section{Stage 1: Mathematical Foundations}

\subsection{Purpose}
This stage defines the financial problem and the notation used throughout the
project. The objective is not to introduce a new option-pricing theory. It is
to explain why early exercise changes a familiar Black--Scholes equation into
an inequality problem with an unknown boundary.

\subsection{Method}
Let \(S>0\) denote the current asset price, \(K>0\) the strike, \(T\) the
maturity date, \(t\in[0,T]\) calendar time, \(r\) the continuously compounded
risk-free rate, \(q\) the continuous dividend yield, and \(\sigma>0\) the
volatility. The project studies American puts and American calls. Their
immediate-exercise payoffs are
\[
  \payoff_{\mathrm{put}}(S)=\max(K-S,0),
  \qquad
  \payoff_{\mathrm{call}}(S)=\max(S-K,0).
\]
The American value \(V(S,t)\) is the best discounted expected payoff available
when the holder may choose an exercise time \(\theta\) between \(t\) and \(T\):
\[
  V(S,t)
  =
  \sup_{\theta\in[t,T]}
  \mathbb{E}^{\mathbb{Q}}\!
  \left[
    e^{-r(\theta-t)}\payoff(S_\theta)
    \,\middle|\,S_t=S
  \right].
\]
The expectation is under the risk-neutral measure \(\mathbb{Q}\).

\begin{readerintuition}
A European holder has one allowed exercise date. An American holder repeatedly
compares two actions: exercise now for \(\payoff(S)\), or retain the contract.
The option value must therefore be at least the immediate payoff.
\end{readerintuition}

\subsection{Mathematical Formulation}
\subsubsection{Time to maturity}
The numerical solver uses
\[
  \tau=T-t,
  \qquad
  U(S,\tau)=V(S,T-\tau).
\]
Thus \(\tau=0\) is maturity and \(\tau=T\) is the initial valuation date. This
change converts a terminal-value problem in \(t\) into an initial-value problem
in \(\tau\):
\[
  U(S,0)=\payoff(S).
\]

\begin{commonconfusion}
Increasing \(t\) moves forward in calendar time, but increasing \(\tau\) moves
away from maturity. The code begins with the payoff at \(\tau=0\) and advances
through \(\tau_1,\tau_2,\ldots,\tau_N=T\).
\end{commonconfusion}

\subsubsection{Risk-neutral dynamics and the Black--Scholes operator}
Under the model used in the project,
\[
  dS_t=(r-q)S_t\,dt+\sigma S_t\,dW_t.
\]
In a region where exercise is not optimal, standard replication and
risk-neutral valuation give
\[
  U_\tau=\mathcal{L}U,
\]
where
\[
  \mathcal{L}U
  =
  \frac12\sigma^2S^2U_{SS}
  +(r-q)SU_S-rU.
\]
The three terms have separate meanings:
\begin{itemize}
  \item \(\tfrac12\sigma^2S^2U_{SS}\) is the diffusion or curvature term;
  \item \((r-q)SU_S\) is the risk-neutral drift term;
  \item \(-rU\) is the discount term.
\end{itemize}

\subsubsection{From a PDE to a variational inequality}
Immediate exercise is always available, so
\[
  U(S,\tau)\geq\Phi(S).
\]
If \(U(S,\tau)>\payoff(S)\), retaining the option is more valuable and the
continuation PDE applies:
\[
  U_\tau-\mathcal{L}U=0.
\]
If \(U(S,\tau)=\payoff(S)\), the holder exercises. At such a point the PDE
equality need not hold, but delaying exercise cannot produce a locally better
value. The combined problem is
\[
  U(S,\tau)\geq\payoff(S),
  \qquad
  U_\tau-\mathcal{L}U\geq0,
  \qquad
  \bigl(U-\payoff\bigr)
  \bigl(U_\tau-\mathcal{L}U\bigr)=0.
\]

\begin{readerintuition}
The product condition says that at least one factor is zero. In continuation,
\(U-\payoff>0\), so the PDE residual must be zero. In exercise,
\(U-\payoff=0\), while the residual may be positive. This is why American
pricing is an inequality or complementarity problem rather than only a linear
PDE solve.
\end{readerintuition}

\subsubsection{Exercise region, continuation region, and free boundary}
Define
\[
  \mathcal{E}
  =
  \{(S,\tau):U(S,\tau)=\payoff(S)\},
  \qquad
  \mathcal{C}
  =
  \{(S,\tau):U(S,\tau)>\payoff(S)\}.
\]
The interface between \(\mathcal{E}\) and \(\mathcal{C}\) is the free boundary
\(B(\tau)\). It is ``free'' because its location is not prescribed as an input;
it must be discovered together with the option value.

For a put, exercise generally occupies the low-spot side:
\[
  S\lesssim B_{\mathrm{put}}(\tau)
  \quad\Longrightarrow\quad
  \text{exercise}.
\]
For a dividend-paying call, exercise---when it occurs---generally occupies the
high-spot side:
\[
  S\gtrsim B_{\mathrm{call}}(\tau)
  \quad\Longrightarrow\quad
  \text{exercise}.
\]
A no-dividend American call is retained as a control because, under the
standard assumptions, early exercise is not optimal.

\subsubsection{Continuation premium}
The continuation premium is
\[
  C(S,\tau)=U(S,\tau)-\Phi(S).
\]
The obstacle condition immediately gives
\[
  C(S,\tau)\geq0.
\]
Moreover,
\[
  C=0 \quad\text{in the exercise region},
  \qquad
  C>0 \quad\text{in the continuation region}.
\]

\begin{readerintuition}
\(C\) is the extra value of keeping the option alive instead of exercising
immediately. It is therefore both a financial quantity and a useful numerical
indicator of the exercise boundary.
\end{readerintuition}

\begin{commonconfusion}
Nonnegative continuation premium is only the payoff-obstacle requirement. It
does not by itself guarantee monotonicity, convexity, a correct free boundary,
or accurate Delta and Gamma.
\end{commonconfusion}

\subsection{Implementation Details}
The solver stores spot along one array axis and time to maturity along the
other. Reports use normalized quantities such as \(S/K\), \(U/K\), and \(B/K\)
where useful, but the mathematical obstacle remains \(U\geq\payoff\).
Later neural features include
\[
  x=
  \bigl(
    \log(S/K),\tau/T,r,q,\sigma,T,
    \mathbf{1}_{\mathrm{call}}
  \bigr).
\]

\subsection{Validation and Diagnostics}
Before discretization, the notation is checked for the direction of time, the
sign of \(U_\tau-\mathcal{L}U\), put/call payoff orientation, and the units of
\(\Delta=U_S\) and \(\Gamma=U_{SS}\). A sign error can yield stable-looking
arrays while solving the wrong problem.

\subsection{Results}
Stage 1 produces a consistent mathematical contract. It introduces no new
numerical results and makes no model-performance claim.

\subsection{Failure Modes and Limitations}
The formulation assumes the Black--Scholes diffusion with constant parameters.
It does not include jumps, stochastic volatility, transaction costs, or market
calibration. These omissions delimit the project rather than invalidate its
within-model numerical study.

\subsection{Conclusion and Next Step}
The continuous problem is an obstacle-constrained PDE. Stage 2 converts its
spatial operator into a tridiagonal matrix, applies Crank--Nicolson in time, and
solves the resulting linear complementarity problem with PSOR.
